\newpage
\section{The Propers: Detailed Commentary}
\subsection{Festive joy and the request for mercy}
The Entrance's incipit, \textit{Gaudeamus omnes in Domino}, means ``Let us all rejoice in the Lord.'' Its subject is common rejoicing in the saints' celebration, with angelic praise of God's Son. It is a composed antiphon. The chant-book directions sometimes printed beside it are separate musical provisions; they add no appointed psalm to these readings. The joy expressed at the entrance concerns a communion larger than the assembly now present.

The Collect asks for divine mercy with the saints' intercession. Celebration and petition occur together: the congregation praises what God has done in the saints and still needs his favor. That combination becomes intelligible beside 1 John. Children may truly belong to the Father while still awaiting completion. The saints' intercession addresses that continuing need. The incipit \textit{Omnipotens sempiterne Deus} addresses almighty and eternal God; the complete prayer is located in the appointed-text section below rather than supplied in a new translation.

\subsection{Revelation: who can stand before God?}
The First Reading occurs between the sixth and seventh seals. At the end of chapter 6, people hide from the wrath of the enthroned one and the Lamb, and ask who can stand. Chapter 7 depicts servants sealed before destruction and then a multitude standing before the throne. In chapter 5 the Lamb has appeared as slain and has received the book. These scenes explain why the multitude attributes salvation to God and the Lamb. Its security comes through the one whose death is already visible in the vision.\footnote{Revelation 5:6--10; 6:12--17; 7:2--4,9--14, Douay--Rheims/Challoner.}

The appointment retains the number of the sealed but passes over the twelve tribal counts in verses 5--8. The hearer therefore encounters a definite total and then a crowd no one can number. Bede explains the number through twelve multiplied by twelve and by a thousand. He associates twelve with the apostles and with faith in the Trinity throughout the fourfold world. His conclusion is the fullness of the Church, descended from the patriarchs by natural lineage or imitation of their faith. The arithmetic is symbolic exegesis of an apocalyptic vision.\footnote{Bede, \work{Explanation of the Apocalypse}, trans. Edward Marshall (Oxford and London: James Parker, 1878), I, on 7:4, pp. 45--46.}

At verse 9 Bede permits two explanations of the transition. The vision may resume the account of victorious glory after an intervening recapitulation. It may also proceed from Israel, to whom the Gospel was first preached, to the salvation of the Gentiles. The second possibility distinguishes the groups at that point in the narrative. Both preserve the breadth of the vision; his discussion supplies no ceiling on the number who can be saved.\footnote{Bede, \work{Explanation}, I, on 7:9, p. 52.}

The elder's question about the white-robed crowd makes their identity depend on what they have come through and how they were cleansed. Blood ordinarily stains a garment. Here the Lamb's blood makes it white. Bede explicitly includes the whole Church in this cleansing, beyond those who died as martyrs. He connects their standing together in God's service with faithful confession amid adversity. The received commentary of Victorinus similarly understands the robes through baptism in the Lamb's blood and the preservation of grace received. Its transmitted text has a difficult recension history; that baptismal explanation is cited from the received commentary, without assigning every phrase to an unaltered ancient original.\footnote{Bede, \work{Explanation}, I, on 7:14, p. 53; received Victorinus, \work{Commentary on the Apocalypse}, seventh chapter, comment 9, ANF VII.}

\subsection{Psalm 24: the Creator and those who seek his face}
The psalm begins with God's ownership of the world and its inhabitants. It then asks who can ascend his mountain and remain in his holy place. The answer concerns hands, heart and the soul's attachment to vanity. The worshipper receives blessing and mercy, and belongs to the generation seeking God's face. The complete psalm continues with the entrance of the King of glory; the appointed strophes end before that dialogue.\footnote{Psalm 23:1--10 in Vulgate/Douay--Rheims numbering, corresponding to Hebrew Psalm 24.}

Augustine reads the whole psalm in relation to Christ's resurrection and glorification. The nations believing in the glorified Lord become his Church. He takes the mountain as the height of God's righteousness, innocent hands as guiltlessness in deed, and the pure heart as purity in thought. Refusing vanity concerns a soul that longs for lasting eternity instead of treating itself as another passing possession. His account combines moral action with desire for an enduring good.\footnote{Augustine, \work{Expositions on the Psalms}, Psalm 23, §§1--6; NPNF I.8 heads it ``Psalm XXIV,'' with the explicit Latin number XXIII.}

Aquinas distinguishes approaching the mountain from standing at the destination. He then distinguishes exterior innocence from interior purity and directly cites Matthew's promise that the clean of heart will see God. His explanation of blessing and mercy has two forms. They can be the promised reward. They can also explain how a person becomes innocent enough to approach: God's mercy makes such access possible. That second explanation matters when the psalm is read beside John's command to purify oneself.\footnote{Thomas Aquinas, \work{Super Psalmos} 23, nn. 2--4, Corpus Thomisticum, paragraphs 86989--86991.}

The response is appointed as an adaptation, ``cf.'' verse 6. The historical study Bible gives a third-person description of the seekers, and does not supply the approved adapted refrain. Its wording can explain the response's scriptural basis; it cannot be recast as an invented text for the assembly to sing.

\subsection{First John: children now, likeness still to come}
The epistle places the Father's love before the believer's work. The children already possess their identity; the world fails to recognize them because it failed to recognize him. Yet the author distinguishes that present identity from what will be manifested. Future likeness and the vision of God explain why hope has a purifying effect now. The next verses address sin and righteousness, and the wider chapter proceeds to love expressed in deeds. Adoption has moral consequences.\footnote{1 John 2:28--3:18; the appointed passage is 3:1--3.}

Augustine's Tractate IV follows the passage from abiding in Christ into these assertions about sonship. He asks what remains to be hoped for if believers are already God's children. His answer distinguishes seeing Christ's humanity at judgment from the promised sight of the divine Word. The wicked too can see the judge in his human form; the promise to the pure concerns a vision beyond that encounter. He then explains desire as preparation for receiving what is promised. The actual reasoning about desire and willing purification is developed in the first interpretation below.\footnote{Augustine, \work{Tractates on the First Epistle of John} IV.4--7, NPNF I.7.}

\subsection{Christ's invitation before the Beatitudes}
Matthew 11:28 follows Jesus' thanksgiving for the revelation given to little ones and his assertion of the Son's knowledge of the Father. Its invitation addresses the burdened. The following verses speak of learning meekness and humility and taking his yoke. Only verse 28 is appointed as the acclamation. The context explains what the offered rest involves.\footnote{Matthew 11:25--30.}

Chrysostom asks how a yoke can offer rest when Jesus elsewhere demands a narrow way and costly obedience. He answers through meekness and humility: the disciple learns from Christ and ceases to bear the burdens of pride and sin in the same way. He retains both labor and gift. The invitation begins with the burden people already carry; Christ calls them toward salvation. Chrysostom explicitly relates this humility to the beginning of Jesus' earlier teaching, the Beatitudes.\footnote{John Chrysostom, \work{Homilies on Matthew} XXXVIII.3, NPNF I.10.}

\subsection{The Gospel: blessing, conduct and promised fulfillment}
Jesus sees the crowds, ascends a mountain, sits and teaches the disciples who approach him. Matthew has just described his Galilean teaching and healing. The discourse later extends through commands about reconciliation, prayer and obedience. Its concluding comparison of houses concerns hearing these words and doing them. Augustine uses that conclusion to explain why the Sermon describes the Christian life in its fullness.\footnote{Matthew 4:23--5:2; 7:24--29; Augustine, \work{Our Lord's Sermon on the Mount} I.1.1--3, NPNF I.6.}

The promises vary: kingdom, inheritance, consolation, satisfaction, mercy, vision and sonship. Augustine understands their variety through a course of conversion beginning with humility. Chrysostom stresses that the initial blessing concerns chosen humility. Low social position can coexist with pride, and disciples honored for their works will still need humility. His explanation makes the invitation available across social ranks while requiring a real disposition of the will.\footnote{Augustine, \work{Sermon on the Mount} I.2--4, §§4--12; Chrysostom, \work{Matthew} XV.2--3.}

Augustine understands the earth as a stable, eternal inheritance. Chrysostom includes a present, sensible reward as well as future gifts: violent striving can lose property and life, while meekness can preserve them. Aquinas reports Chrysostom and prefers a future interpretation.\footnote{Augustine, \work{Sermon on the Mount} I.2.4; Chrysostom, \work{Matthew} XV.5; Aquinas, \work{Super Matthaeum}, cap. 5, lectio 2, paragraph 87340.}

The Gospel ends after the promise of heavenly reward in 12a. The following reference to earlier prophets belongs to the complete biblical context and Chrysostom's commentary, rather than to the text appointed for proclamation. Likewise the study Bible and Augustine place meekness before mourning; the U.S. proclaimed Gospel places mourning first. Augustine's explanation of successive stages follows the sequence of his own text.

\subsection{The saints' concern at the Eucharist}
The Prayer over the Offerings concerns the saints' continuing concern for the salvation of those still on the way. Their completed life is thus compatible with help for others. The prayer is addressed to God, and their intercession is understood within that dependence. Its subject belongs beside the Collect's request for mercy, although the two prayers occupy different ritual moments.

The proper Preface names Jerusalem as mother and speaks of the saints' glory, the pilgrimage toward it, and their help and example. These are personal and communal descriptions. The living assembly can learn from those whose journey is complete and receive their intercession. The Preface's Jerusalem language is part of a composed prayer; it does not appoint an additional reading from Revelation 21 or Hebrews 12. Its exact English remains available at the publisher locator below.

At Communion the antiphon selects Matthew 5:8--10: purity, peacemaking and persecution for justice. Vision, sonship and kingdom are consequently heard together at this point. Purity concerns the recipient's heart, peacemaking the work done among persons, and persecution the cost fidelity may have. The selection keeps these demands beside one another. The two interpretations explain their relation in different ways.

The Prayer after Communion concerns charity's perfection and arrival at the heavenly banquet. Receiving Communion is followed by asking for further grace. The future banquet is hoped for, while the sacrament is received now. The prayer therefore permits an account of real participation that still needs perseverance and growth. If the solemn blessing is selected, its petitions concern intercession, holy example, service and final joy together. Its application to service is pertinent to both interpretations, but it is an optional conclusion.\footnote{Missal prayer subjects and the permitted blessing are controlled by the All Saints sources identified in References; withheld English bodies are not translated here.}
